
\subsection{Objetivo}

\subsubsection{General}
Desarrollar un sistema móvil de visión por computadora, basado en aprendizaje automático, para identificar automáticamente la presencia de plagas de insectos descortezadores y clasificar el estado sanitario del arbolado en un bosque en específico de la Ciudad de México. El sistema utilizará imágenes recolectadas en campo y permitirá la detección temprana de focos de infestación para fomentar su atención.

\subsubsection{Específicos}

\begin{enumerate}
    \item Identificar el bosque en la Ciudad de México más apto para la construcción del \textit{dataset}.
    \item Construir y etiquetar datos tomados directamente en campo de imágenes del arbolado, identificando casos sanos e infectados.
    \item Diseñar una arquitectura en la nube que permita almacenar, procesar y entrenar al modelo de manera eficiente.
    \item Entrenar un modelo de visión por computadora capaz de reconocer síntomas de la infestación por descortezadores en los árboles de un bosque en específico de la Ciudad de México.
    \item Desarrollar una aplicación móvil para la captura de imágenes y seguimiento de los resultados.
\end{enumerate}